% section_1_intro.tex — Introduction
% Drafted per: docs/Draft/THESIS_SKELETON.md v5.1 §I (post-capex-drop 2026-04-22):
%   Motivation, gap, approach, empirical setting + findings, contributions, roadmap.
% Citation discipline: only committed-bib entries (12); CEO-as-face rationale forwarded to §2.3.

\section{Introduction}
\label{sec:intro}

\subsection{Motivation}

Corporate earnings calls offer a quarterly window into managerial communication that is higher in frequency than the annual 10-K disclosure cycle and richer in source-identifiable speaker attribution: a single call admits separate prepared-remarks and Q\&A segments and identifies who is speaking (CEO, CFO, other manager, analyst). The central question we ask is whether uncertainty conveyed by the chief executive officer specifically, in either segment of the call, predicts the firm's contemporaneous and one-quarter-lead cash holdings---the precautionary response formalized by \citeA{opler1999}: firms facing greater uncertainty about future cash flows hold more cash as a liquidity buffer against future funding shocks. \citeA{minton2001} document a cross-sectional cluster of low-leverage, high-cash firms whose pattern they label \emph{financial conservatism}; the present paper inherits the noun (\emph{Financing Conservatism}) for its title, but the empirical test targets the cash margin of that posture, leaving the leverage margin to future work.

\subsection{Gap}

The finance-specific uncertainty word list of \citeA{loughran2011} has been deployed extensively at the annual 10-K disclosure frequency, and \citeA{dzielinski2021} and \citeA{hassan2020} have extended textual uncertainty measurement to the higher-frequency earnings-call corpus at firm-quarter resolution. To our knowledge, however, no prior work has systematically applied the CEO-specific speaker partition---enabled by the structured speaker-metadata fields in the Capital~IQ vendor---to test the precautionary cash response of \citeA{opler1999} at firm-quarter frequency, nor has prior work jointly examined the firm's own cash-financing margin and the outside-world reaction margin (post-call quoted bid-ask spread) of CEO speech uncertainty. The CEO partition matters substantively in our setting: the empirical decomposition reported in Appendix~\ref{app:partition} shows that pooling CEO speech with non-CEO manager speech, as in the all-manager aggregation convention of \citeA{bushee2018}, masks the concentration of the cash-holdings response in the CEO Q\&A segment.

\subsection{Approach}

The speech-uncertainty regressor is constructed by applying the \citeA{loughran2011} uncertainty word list to earnings-call transcripts in the segment-split convention of \citeA{dzielinski2021}, who document that a single call carries systematically different uncertainty content in its prepared-remarks and Q\&A segments. The aggregation architecture follows the all-manager pooling of \citeA{bushee2018}; we restrict that pooling to the CEO-tagged subset of speakers using the Capital~IQ structured speaker-metadata fields. The Q\&A regressor $\text{UncAnsCEO}$ counts uncertainty words in CEO Q\&A utterances divided by total CEO Q\&A word count, and the Presentation regressor $\text{UncPreCEO}$ is defined analogously on the prepared-presentation segment. The rationale for the CEO-specific partition---rather than aggregating over the full manager pool---is documented in \S\ref{sec:framework:measurement}; in brief, the choice rests on the upper-echelons role of the chief executive officer as principal external-facing communicator of the firm and on the empirical concentration of the cash-holdings response in the CEO Q\&A segment.

\subsection{Empirical Setting and Findings}

The sample comprises $112{,}968$ earnings-call transcripts covering $2{,}429$ unique firms over fiscal years 2002--2018, sourced from S\&P Capital~IQ and merged with Compustat quarterly fundamentals and CRSP daily security data; financial firms (SIC 6000--6999) and regulated utilities (SIC 4900--4999) are excluded throughout, following \citeA{bates2009}. The main empirical analyses test two hypotheses pre-committed in \S\ref{sec:framework:precommit}: HC (the cash response, \S\ref{sec:main:hc}) and HFC (the financial-constraint moderation of HC, \S\ref{sec:main:hfc}). For HC, UncAnsCEO is positive and statistically significant in all twelve specifications of suite H1 across a four-step fixed-effects ladder, while the Presentation channel UncPreCEO is null across the same panel. For HFC, estimated on the 2002--2016 sub-sample where Compustat credit-rating coverage is available, the UncAnsCEO\_c~$\times$~Unrated interaction is positive and significant in six of eight interaction cells, with all four lead-dependent-variable cells reaching $p<0.10$; the precautionary cash response therefore concentrates in the credit-constrained Unrated firms of \citeA{faulkender2006}. Two additional sets of analyses follow. The driver regressions (\S\ref{sec:additional:drivers}) test whether the CEO speech-uncertainty metric itself responds to independently measured exogenous uncertainty: the firm-level political-risk index of \citeA{hassan2020} loads both CEO channels at high precision, and the news-based macro indices of \citeA{baker2016} and \citeA{davis2016} load the Presentation channel preferentially. The outside-world reaction regressions (\S\ref{sec:additional:reaction}) test whether liquidity providers price CEO speech uncertainty into post-call quoted spreads: the 26-day post-call closing-quote relative bid-ask spread (Spread\textsubscript{25D}) widens contemporaneously with the Presentation channel and at a one-quarter lag with the Q\&A channel. The Q\&A-versus-Presentation channel asymmetries documented across the cash, driver, and spread analyses are reported descriptively, consistent with the pre-commitment of \S\ref{sec:framework:precommit} that channel-specific asymmetry is reported ex post rather than asserted ex ante.

\subsection{Contributions}

The paper makes three contributions to the literature on managerial uncertainty disclosure and corporate financing.
\begin{enumerate}
\item \textbf{Systematic CEO-partitioned application to firm-quarter financing.} The framework of \citeA{dzielinski2021} establishes that conference-call transcripts admit a CEO-specific uncertainty metric at the individual-speaker level, and the architecture of \citeA{bushee2018} establishes the canonical pooled-manager aggregation. The principal measurement-side contribution of this paper is the systematic application of the CEO-specific partition---rather than the pooled-manager aggregate---to firm-quarter cash-financing outcomes (HC, HFC) and to firm-quarter external-reaction outcomes (\S\ref{sec:additional:reaction}). The CEO-partitioned regressor isolates a cash-response concentration that the pooled-manager aggregate obscures, as documented in Appendix~\ref{app:partition}.
\item \textbf{Multi-aggregation-level construct validation of the CEO speech-uncertainty metric.} The driver regressions of \S\ref{sec:additional:drivers} provide construct-validation evidence that the CEO-partitioned measure carries economically meaningful uncertainty content at firm, US-macro, and global-macro aggregation levels, responding positively to the political-risk index of \citeA{hassan2020} and the news-based policy-uncertainty indices of \citeA{baker2016} and \citeA{davis2016}. The validation strengthens the substantive interpretation of HC and HFC as identifying an uncertainty channel of CEO speech rather than a stylistic artifact of the wordlist or the segment-split.
\item \textbf{Outside-world reaction in market liquidity.} The post-call spread regressions of \S\ref{sec:additional:reaction} document that liquidity providers price both CEO speech segments into post-call closing-quote relative bid-ask spreads, with the Presentation channel registering contemporaneously within the call quarter and the Q\&A channel registering with a one-quarter lag. The outside-world reaction margin is theoretically distinct from the firm-side cash response of \S\ref{sec:main}, and the two findings are presented as descriptively complementary rather than as a tested causal sequence linking firm financing to market microstructure.
\end{enumerate}

\subsection{Roadmap}

The remainder of the paper is organized as follows. Section~\ref{sec:framework} states the conceptual framework (\S\ref{sec:framework:conceptual}), the speech-uncertainty measurement architecture (\S\ref{sec:framework:measurement}), the precautionary-motive and financing-conservatism literature anchors for HC and HFC (\S\ref{sec:framework:precautionary}), and the empirical design including the fixed-effects ladder and clustering conventions (\S\ref{sec:framework:design}). Section~\ref{sec:main} reports the data, sample, variables, and specification (\S\ref{sec:main:setup}), the HC results (\S\ref{sec:main:hc}), and the HFC results (\S\ref{sec:main:hfc}). Section~\ref{sec:additional} reports the construct-validation driver regressions (\S\ref{sec:additional:drivers}) and the outside-world spread reaction (\S\ref{sec:additional:reaction}). Section~\ref{sec:conclusion} synthesizes the channel asymmetries across the three empirical sections, discloses limitations, and outlines five extensions for future research.
